\section*{Резюме}

Инжекционните атаки остават сред най-разпространените уязвимости в софтуера и са
средство за редица документирани киберпрестъпления. Настоящата дипломна работа
изследва дали една абстракция, конструираща заявките към хранилища за данни по
време на компилация, може да елиминира цяло семейство от класове инжекции при
хетерогенни хранилища, и къде точно спира тази гаранция. Като платформа се
използва библиотека за обектно-релационно картографиране на C++23, разработена в
предходна дипломна работа, която работи със седем хранилища: PostgreSQL, MySQL,
SQLite, Cassandra, MongoDB, Neo4j и Redis.

Въз основа на определението на инжекцията като промяна на синтактичната
структура на заявката е формулиран аргумент: когато имената на таблици и колони и
формата на заявката са аргументи на шаблони, а стойностите достигат заявката само
чрез свързани типизирани заместители, недоверен вход не може да промени
структурата. Аргументът е проверен с три независими метода: 26 отрицателни
компилационни теста и 5 контролни случая, с проверка на диагностиката при GCC 14.2
и Clang 18.1; диференциален тест, при който всеки приложим полезен товар от корпус
от 66 публикувани полезни товара от 11 източника се подава към съзнателно уязвима
заявка и към библиотеката; и анализ на изходния код на всеки конектор.

Анализът установи, че гаранцията се постига с четири различни механизма, и откри
пропуск в нея: стойност, записана направо в заявката вместо чрез заместител, се
приемаше, но не се свързваше, и заявката тихо връщаше грешни резултати. Пропускът
е затворен с проверка по време на компилация във всички седем входни точки, при
запазено успешно изпълнение на всички съществуващи тестове. При SQLite
наивната заявка беше компрометирана от 4 от 22 приложими полезни товара,
включително с извличане на съхранена идентификационна информация, а библиотеката
от нито един. Тестът беше изпълнен и срещу работещи сървъри на PostgreSQL, MySQL,
MongoDB, Redis, Cassandra и Neo4j в система за непрекъсната интеграция. Общо за
седемте хранилища през библиотеката не е имала ефект нито една от 98-те
двойки на полезен товар и хранилище, съставени от всички 66 полезни товара.
Наивният път е компрометиран от по 2 полезни товара при PostgreSQL, MySQL,
MongoDB, Redis и Neo4j. При Cassandra той не е компрометиран от нито един и
сравнението е неубедително.

Библиотеката беше разширена в две посоки, в които гаранцията не е очевидна. При
условия \code{IN} и \code{NOT IN} вход по време на изпълнение определя дължината
на списък, но текстът на заявката получава само броя на маркерите. Миграциите
работят с имена, прочетени обратно от самото хранилище, затова там гаранцията е
проверка по време на изпълнение, а не компилация. Тя е измерена и при инжекция от
втори ред: наивният DDL е компрометиран при 14 от 368 изпълнения,
всички от съставени входове, а библиотеката при нито едно.

Приносите на работата са аргумент за предотвратяване на инжекциите чрез
конструиране по време на компилация, разширен извън SQL; възпроизводим набор за
оценка с корпус с посочени източници; открит и затворен пропуск в гаранцията; и
определена граница на гаранцията при списъци и операции по схемата.

\bigskip
\noindent\textbf{Ключови думи:} инжекционни атаки, предотвратяване на
киберпрестъпления, сигурност по проектиране, C++, програмиране по време на
компилация, хранилища за данни.

\section*{Abstract}

Injection attacks remain among the most common software vulnerabilities and are
the means of a number of documented cybercrimes. This thesis examines whether an
abstraction that constructs queries to data stores at compile time can eliminate
a whole family of injection classes across heterogeneous stores, and where
exactly that guarantee stops. The platform is a C++23 object-relational mapping
library, developed in an earlier thesis, that works with seven stores:
PostgreSQL, MySQL, SQLite, Cassandra, MongoDB, Neo4j and Redis.

Building on the definition of injection as a change to the syntactic structure
of a query, the thesis argues that when table names, column names and the shape
of the query are template arguments, and values reach the query only through
bound, typed placeholders, untrusted input cannot change that structure. The
argument is tested by three independent methods: 26 negative compilation tests
and 5 control cases, with diagnostic matching under GCC 14.2 and Clang 18.1; a
differential test that sends every applicable payload from a corpus of 66
published payloads from 11 sources both to a deliberately vulnerable query and to
the library; and an analysis of the source of every connector.

The analysis found that the guarantee is achieved by four different mechanisms,
and found a gap in it: a value written directly into a query instead of through a
placeholder was accepted but never bound, and the query silently returned wrong
results. The gap is closed by a compile-time check at all seven entry points,
with all existing tests still passing. On SQLite the vulnerable query
was compromised by 4 of 22 applicable payloads, including the theft of a stored
credential, and the library by none. The test was then run against live servers
of PostgreSQL, MySQL, MongoDB, Redis, Cassandra and Neo4j in a continuous
integration system. Across the seven stores, none of the 98 pairs of payload
and store, covering all 66 payloads, had an effect through the library. The naive
path was compromised by 2 payloads each on PostgreSQL, MySQL, MongoDB, Redis and
Neo4j. On Cassandra it was never compromised, so the comparison is inconclusive.

The library was extended in two directions where the guarantee is not obvious. For
\code{IN} and \code{NOT IN}, run-time input sets the length of a list, but the query
text receives only the number of markers. Migrations work with names read back from
the store itself, so there the guarantee is a run-time check and not a compile-time
one. It was measured, including second order injection: the naive DDL was compromised
in 14 of 368 runs, all by constructed inputs, and the library in none.

The contributions are an argument for preventing injection by compile-time
construction, extended beyond SQL; a reproducible evaluation suite with a sourced
corpus; a gap in the guarantee, found and closed; and the boundary of the guarantee
for lists and schema operations, located and measured.

\bigskip
\noindent\textbf{Keywords:} injection attacks, cybercrime prevention, secure by
design, C++, compile-time programming, data stores.
